\documentclass[11pt,a4paper]{article}

\usepackage[T1]{fontenc}
\usepackage[utf8]{inputenc}
\usepackage{libertinus}
\usepackage{microtype}
\usepackage[a4paper,margin=21mm,headheight=15pt]{geometry}
\usepackage{amsmath,amssymb,mathtools}
\usepackage{booktabs,longtable,tabularx,array}
\usepackage{enumitem}
\usepackage[most]{tcolorbox}
\usepackage{xcolor}
\usepackage{fancyhdr}
\usepackage[numbers,sort&compress]{natbib}
\usepackage{hyperref}

\definecolor{navy}{HTML}{163A5F}
\definecolor{teal}{HTML}{087E8B}
\definecolor{green}{HTML}{2A7F62}
\definecolor{pale}{HTML}{EEF5F7}
\definecolor{warm}{HTML}{FFF6E8}
\definecolor{redpale}{HTML}{FFF0F0}
\definecolor{muted}{HTML}{4B5563}
\hypersetup{colorlinks=true,linkcolor=navy,citecolor=teal,urlcolor=teal,
  pdfauthor={DMD initial validation planning},
  pdftitle={Initial validation plan for DMD on deconvolved calcium-event data}}
\pagestyle{fancy}
\fancyhf{}
\lhead{DMD initial validation plan}
\rhead{Version 1.1 --- 16 July 2026}
\cfoot{\thepage}
\setlength{\parindent}{0pt}
\setlength{\parskip}{0.52em}
\setlist{nosep,leftmargin=1.6em}
\renewcommand{\arraystretch}{1.16}
\newcolumntype{Y}{>{\raggedright\arraybackslash}X}
\newcolumntype{P}[1]{>{\raggedright\arraybackslash}p{#1}}

\newtcolorbox{decisionbox}[1][]{breakable,colback=pale,colframe=teal,
  boxrule=0.75pt,arc=2pt,left=7pt,right=7pt,top=6pt,bottom=6pt,
  title=#1,fonttitle=\bfseries}
\newtcolorbox{gatebox}[1][]{breakable,colback=green!5,colframe=green,
  boxrule=0.75pt,arc=2pt,left=7pt,right=7pt,top=6pt,bottom=6pt,
  title=#1,fonttitle=\bfseries}
\newtcolorbox{riskbox}[1][]{breakable,colback=warm,colframe=orange!70!black,
  boxrule=0.75pt,arc=2pt,left=7pt,right=7pt,top=6pt,bottom=6pt,
  title=#1,fonttitle=\bfseries}
\newtcolorbox{stopbox}[1][]{breakable,colback=redpale,colframe=red!65!black,
  boxrule=0.75pt,arc=2pt,left=7pt,right=7pt,top=6pt,bottom=6pt,
  title=#1,fonttitle=\bfseries}

\newcommand{\T}{^{\mathsf T}}
\newcommand{\Hh}{^{\mathsf H}}
\newcommand{\R}{\mathbb{R}}
\DeclareMathOperator{\tr}{tr}

\begin{document}
\hypersetup{pageanchor=false}

\begin{titlepage}
  \centering
  \vspace*{1.6cm}
  {\color{navy}\rule{\textwidth}{1.2pt}}\\[1.0cm]
  {\Huge\bfseries Does DMD recover reproducible dynamics\\[0.2em]
  from deconvolved calcium-event data?\par}
  \vspace{0.55cm}
  {\Large Initial validation plan before the full brain-state analysis\par}
  \vspace{0.9cm}
  {\color{navy}\rule{\textwidth}{0.6pt}}\\[0.9cm]
  \begin{tcolorbox}[colback=pale,colframe=teal,width=0.92\textwidth,arc=2pt]
  \centering
  \textbf{Planning checkpoint only.} No DMD test or biological comparison has
  been run for this document. Implementation will begin only after this plan is
  reviewed. The workflow stops again after the initial validation report so that
  its results can be inspected before any all-recording analysis or mode tracking.
  \end{tcolorbox}
  \vfill
  \begin{tabular}{rl}
    Project: & \texttt{ca\_imaging\_network\_analysis/DMD}\\
    Version: & 1.1\\
    Date: & 16 July 2026\\
    Parent plan: & \texttt{02\_dmd\_brain\_state\_analysis\_plan.pdf}
  \end{tabular}
  \vfill
\end{titlepage}

\hypersetup{pageanchor=true}
\pagenumbering{roman}
\tableofcontents
\clearpage
\pagenumbering{arabic}

\section{Decision summary}

\begin{decisionbox}[Narrow decision to be made]
The pilot asks whether at least one neuron-resolved observation and preprocessing
arm yields a low-rank linear propagator that is (i) predictively useful on
contiguous held-out data, (ii) numerically and bootstrap reproducible, (iii)
distinguishable from smoothing and independent-neuron nulls, and (iv) recoverable
under a small known-system simulation matched to empirical sparsity. It does
\emph{not} ask whether awake and unconscious states differ.
\end{decisionbox}

The closest direct precedent is Pachitariu et al., who applied a time-lagged,
ridge-regularized DMD to neuronwise standardized, spike-deconvolved two-photon
data after PCA, and deliberately used a 0.23-s lag to reduce sensitivity to
short-timescale deconvolution artefacts \citep{pachitariu2026critical}. We will
reproduce the logic of that analysis, but not copy parameters that are
incompatible with the present sampling rate and 300-frame target window.

The pilot has three possible conclusions:

\begin{description}[style=nextline,leftmargin=1.7em]
  \item[PASS to expanded validation.] At least one arm passes the synthetic,
  predictive, reproducibility, and tracking-resolution gates. It becomes a
  candidate for validation on all recordings; no biological claim follows yet.
  \item[REVISE.] DMD is useful only after changing the signal, lag, rank, or
  window length. The parent plan is prospectively amended before further work.
  \item[STOP.] No arm is predictive and reproducible beyond the controls, or
  apparent success is reproduced by the smoothing/null pipeline. Intensive DMD
  and Grassmann tracking are not pursued with these data.
\end{description}

\section{Questions and non-goals}

\subsection{Pilot questions}

\begin{enumerate}
  \item Can a ridge-DMD model predict held-out population activity better than
  persistence at physical horizons near 1 and 2 s?
  \item Does the full population operator add information beyond independent
  per-PC autoregressions, or is predictability explained by marginal smoothing
  and calcium autocorrelation?
  \item Are the leading eigenvalues and invariant subspace reproducible under
  moving-block resampling and deterministic neuron subsampling?
  \item Can the estimator recover known real-relaxation versus rotational
  dynamics after sparse-event and calcium-like observation transformations?
  \item Which observation arm---native deconvolution, binned events, stored
  smoothed events, or $\Delta F/F$---is empirically defensible for the next phase?
\end{enumerate}

\subsection{Explicit non-goals}

\begin{stopbox}[Work excluded from this checkpoint]
The pilot will not estimate awake--sleep or awake--anaesthesia effect sizes;
train a state classifier; use state separation to choose preprocessing; implement
SGOT or another Grassmann tracker; claim individual biological mode identity;
analyze transitions; or run cohort-level inference. In-sample reconstruction is
diagnostic and cannot by itself constitute validation.
\end{stopbox}

\section{What is copied from Pachitariu and what is revised}

Pachitariu et al. standardized each neuron, projected activity into as many as
1,000 PCs, and fitted a ridge map from $X_t$ to $X_{t+\Delta t}$, using a ridge
penalty of 0.01 for two-photon data and $\Delta t\simeq0.23$ s
\citep{pachitariu2026critical}. They compared discrete eigenvalue distributions
with matched simulations rather than relying on a noisy inversion to a
continuous generator. The audited reference is the public
\href{https://github.com/MouseLand/critical_init/tree/v1.0}{MouseLand/critical\_init
release v1.0}; its version identifier will be recorded in the pilot manifest.

\begin{riskbox}[Reference-code audit and declared departures]
Release v1.0 centers each neuron and divides by its root-mean-square deviation
plus $10^{-3}$; this is nearly, but not exactly, a textbook sample $z$-score.
It applies an unnormalized Gram-matrix ridge of 0.01, so that number is not
portable across sample counts or ranks. Its rotation code filters on
$|\lambda|>0.25$, whereas the associated figure caption describes a real-part
threshold. The smoke test will reproduce the code definition; the report will
also show the caption definition and label the discrepancy. The local-window
implementation will use an SVD/method-of-snapshots solve rather than explicitly
squaring the covariance condition number. These are recorded adaptations, not
silent changes.
\end{riskbox}

\begin{decisionbox}[Prospective implementation amendment (version 1.1)]
Before any empirical DMD fit was inspected, the implementation audit fixed four
under-specified points. First, a $k=4$ projector inside a $q=4$ model is the
identity and cannot measure mode stability; therefore only $q\geq8$ models are
eligible for the $k=4$ subspace and tracking gates. Lower ranks remain eligible
for prediction but can support only a conditional outcome. Second, ``maximally
time-spread'' now means the global maximum of the minimum separation between
integer window starts, with the lexicographically earliest tied solution. Third,
development validation now mirrors deployment by fitting the first 80\% and
scoring the last 20\% within each complete development window. Fourth, the
approved raw-label windows remain primary, while S$z$/S$c$ receive a seven-frame
edge-trimmed sensitivity because the stored smoother can carry support across a
nearby label transition. These corrections were frozen in the machine-readable
configuration before examining a DMD result.
\end{decisionbox}

\begin{longtable}{@{}P{0.24\textwidth}P{0.31\textwidth}P{0.37\textwidth}@{}}
\toprule
Element & Pachitariu-aligned choice & Required adaptation here\\
\midrule
\endfirsthead
\toprule
Element & Pachitariu-aligned choice & Required adaptation here\\
\midrule
\endhead
Observation & Per-neuron spike-deconvolved activity & This becomes the anchor,
not a presumed winner; stored smoothed events, binned events, and $\Delta F/F$
are observation-model controls.\\
Scaling & Per-neuron centering and RMS scaling before PCA & Reproduce the
release-v1.0 RMS-plus-$10^{-3}$ rule, estimating both quantities from development
data only. Never standardize each 300-frame window independently.\\
Spatial reduction & PCA to 1,000 dimensions & A near-literal 1,000-neuron
sanity check is possible on \texttt{example\_data}; a 300-frame window cannot
identify 1,000 dimensions, so local fits use low fixed ranks.\\
Operator & Ridge regression in PC coordinates & Retain ridge DMD as the main
estimator. Unregularized DMD is only a numerical benchmark.\\
Lag & 0.23 s to reduce deconvolution artefacts & The closest lag is two frames,
$2/7.65=0.261$ s. Four- and eight-frame lags test 0.523 and 1.046 s; one frame is
reported only as an artefact-sensitive diagnostic.\\
Ridge penalty & Unnormalized $\alpha=0.01$ for two-photon data & Retain the raw
0.01 fit as the precedent anchor, but select from a dimensionless, scale-relative
training-only grid because effective shrinkage changes with dimension, sample
count, and feature scale.\\
Fit duration & Long recordings at 22 Hz & Include a long-block diagnostic, then
separately test the intended 300-frame (39.2-s) deployment window and a
1,500-frame (196.1-s) rescue window.\\
Endpoint & Eigenvalue complexity/rotations & Report this endpoint, but require
forecast and subspace stability as additional feasibility evidence.\\
\bottomrule
\end{longtable}

This is a faithful methodological adaptation rather than a literal replication.
In particular, copying 1,000 PCs into 300 samples would guarantee an
underdetermined local operator, and copying the absolute ridge penalty without
auditing scaling would not reproduce the same regularization strength.

\section{Development data and leakage control}

\subsection{Files and fixed selection rule}

The smoke test uses \texttt{data/raw/example\_data.mat}, which contains 1,000
neurons. Because that example is derived from the first sleep dataset, empirical
evidence instead uses a separate sleep recording with adequate windows in both
labels and the smallest-ROI anaesthesia recording, selected before any DMD result:

\begin{center}
\begin{tabularx}{\textwidth}{@{}lrrY@{}}
\toprule
File & ROIs & Frames & Pilot role\\
\midrule
\texttt{example\_data.mat} & 1,000 & 9,350 & Code/numerical smoke test and
near-literal high-dimensional Pachitariu configuration; not biological evidence.\\
\texttt{mouse02\_sleep.mat} & 6,574 & 18,700 & Six disjoint awake and six
disjoint NREM window candidates, using raw labels rather than
\texttt{used\_frame}.\\
\texttt{mouse07\_ane.mat} & 4,337 & 19,000 & Six disjoint awake and six
disjoint anaesthesia window candidates.\\
\bottomrule
\end{tabularx}
\end{center}

This keeps the first computation small while including both paradigms. The next
phase, if approved, repeats the locked configuration on all ten recordings.

\subsection{Window selection and label masking}

Within each recording--state cell, choose six non-overlapping 300-frame windows
that have one constant raw dataset label and remain inside one acquisition
segment. Globally maximize the minimum separation between their integer start
indices; among tied solutions, use the lexicographically earliest sequence. The project-specific
\texttt{used\_frame} variable is not an eligibility criterion. Split the six
deterministically into three development and three untouched evaluation windows,
alternating in time. This yields 12 untouched empirical evaluation windows.

State labels are permitted only to create approximately homogeneous windows and
balance development data. During hyperparameter selection, their names are
replaced by anonymous stratum identifiers, and no between-label performance or
separability criterion is calculated.

For each recording and signal arm, neuron scaling and the common PCA basis are
learned from the six development windows pooled with equal weight across the two
anonymous strata. Hyperparameters use local blocked validation: within each
development window, fit the first 80\% and score the final contiguous 20\%, then
aggregate across the six anonymously. The PCA basis is then frozen for every
evaluation window.
This is essential: window-specific PCA would create artificial subspace motion
and make later mode comparisons uninterpretable.

All neurons with finite, nonzero variance in the development data are eligible;
no neuron is selected using state differences or DMD performance. The pilot will
report the retained count and repeat the winning configuration on deterministic
equal-$N$ neuron subsets as a dimensional sensitivity.

\section{Candidate observation and preprocessing arms}

All arms preserve one row per neuron. No across-neuron averaging, binary
threshold, or new Gaussian smoothing is applied. OASIS amplitudes are not treated
as calibrated spike counts or spikes per second \citep{friedrich2017,kiyooka2026}.

\begin{longtable}{@{}P{0.12\textwidth}P{0.27\textwidth}P{0.23\textwidth}P{0.30\textwidth}@{}}
\toprule
ID & Signal construction & Scaling & Purpose\\
\midrule
\endfirsthead
\toprule
ID & Signal construction & Scaling & Purpose\\
\midrule
\endhead
P & Native \texttt{spike\_deconv} at 7.65 Hz & Train-only, code-matched
centered RMS scaling & Pachitariu anchor; retains temporal resolution but is
highly sparse.\\
B4 & Non-overlapping four-frame sums of \texttt{spike\_deconv}, divided by
0.523 s & Train-only centered RMS scaling & Reduces zero inflation without the
1.96-s stored smoothing. It is an event-mass-rate proxy, not calibrated firing rate.\\
S$z$ & Stored 15-frame/1.96-s \texttt{spike\_smoothed} at 7.65 Hz & Train-only
centered RMS scaling & Tests the signal used by the Kiyooka network analysis
under Pachitariu-like scaling. No second smoothing is added.\\
S$c$ & Stored \texttt{spike\_smoothed} & Recording/development mean-centering,
no variance scaling & Restricted sensitivity matching the current parent plan;
tests whether $z$-scoring creates or erases structure.\\
F & Native $\Delta F/F$ & Train-only centered RMS scaling & Observation-model
control: asks whether continuous calcium carries reproducible dynamics lost by
deconvolution.\\
\bottomrule
\end{longtable}

A nine-frame/1.176-s binned-event arm is deferred from the core screen because a
39.2-s window would contain only about 33 bins. It may be run as a rank-2 stress
test if all higher-resolution event arms fail. If only $\Delta F/F$ succeeds,
the conclusion is to change the primary observation model, not to declare the
deconvolved signal validated. Calcium-aware latent-dynamics methods remain a
later alternative if separate deconvolution and DMD are inadequate
\citep{wei2020,koh2023}.

Because the stored 15-frame smoother has an approximate seven-frame half-width,
S$z$/S$c$ are additionally rescored after trimming seven samples from each edge
of every selected window. This is a declared boundary-support sensitivity; the
untrimmed raw-label window remains the primary arm definition.

\section{Model and minimal hyperparameter grid}

For a fixed arm, common PCA basis, rank $q$, and lag $\ell$, let
$z_t\in\R^q$ be the PC score. Fit
\begin{equation}
  z_{t+\ell}=Bz_t+\varepsilon_t,
  \qquad
  \widehat B=YX\T\left(XX\T+\alpha I\right)^{-1}.
\end{equation}
The eigenvalues of $\widehat B$ are discrete-lag eigenvalues. The pilot will not
convert them to a continuous generator unless a value is well separated from
zero and the complex-log branch is unambiguous.

For deployment fits, define the scale
$s_X=\tr(XX\T)/q$ and write $\alpha=\eta s_X$. Thus $\eta$ is the ridge penalty
relative to the mean eigenvalue of the predictor Gram matrix. The literal
Pachitariu smoke-test fit instead uses the original unnormalized $\alpha=0.01$.

\subsection{Three fit tiers}

\begin{enumerate}
  \item \textbf{Precedent check.} On \texttt{example\_data}, run arm P with
  the release-v1.0 neuronwise RMS scaling, the maximum numerically supported PCA
  dimension up to 1,000, unnormalized $\alpha=0.01$, and a two-frame lag.
  Reproduce the discrete-eigenvalue plane and rotations-per-tenfold-attenuation
  summary of Pachitariu et al. This checks implementation fidelity, not
  biological validity.
  \item \textbf{Deployment check.} In the 300-frame development/evaluation
  windows, use a common recording-level basis and low ranks. Native arms use
  $q\in\{2,4,8,16\}$, additionally capped at no more than one retained
  dimension per ten eligible training snapshot pairs. B4 contains only 75 bins
  and about 60 fitting bins, so it uses $q\in\{2,4\}$.
  \item \textbf{Long-block diagnostic.} In each of the four recording--label
  cells, select the 1,500-frame block centered in the longest eligible
  constant-label bout, without reference to DMD output. Refit at
  $q\in\{4,8,16,32\}$ using blocked training-only selection and a final
  contiguous 20\% evaluation segment. This is a prespecified rescue test that
  distinguishes estimator failure from insufficient 39.2-s support; it is not
  an opportunity to select a favourable biological contrast.
\end{enumerate}

Native-arm lags are $\ell\in\{1,2,4,8\}$ frames, corresponding to 0.131, 0.261,
0.523, and 1.046 s; two frames is the precedent anchor and one frame is diagnostic.
For S$z$/S$c$, a 15-frame (1.961-s) lag additionally tests the scale of the stored
smoothing kernel. B4 uses one-, two-, and four-bin lags.

The raw precedent value $\alpha=0.01$ is always shown. On development data only,
deployment selection uses
\begin{equation*}
  \eta\in\{10^{-4},10^{-3},10^{-2},10^{-1},1\}.
\end{equation*}
Both $\eta$ and the resulting raw $\alpha$ are recorded. Exact DMD ($\eta=0$)
is a conditioning diagnostic, not an eligible winner when it is unstable.

\section{Held-out empirical validation}

\subsection{Blocked prediction}

In each untouched 300-frame evaluation window, estimate the local operator from
the first 80\% (240 frames, or 60 B4 bins) and evaluate on the final contiguous
20\%. No shuffled pair, acquisition break, state boundary, or artificial
bootstrap-block join may enter $X\rightarrow Y$.

Rolling-origin forecasts reinitialize from the observed state and assess the
nearest available horizons to 0.26, 1.05, and 2.09 s. For baseline $b$, define
\begin{equation}
  \operatorname{Skill}_b(h)=
  1-\frac{\sum_t\lVert z_{t+h}-\widehat z^{\mathrm{DMD}}_{t+h}\rVert_2^2}
          {\sum_t\lVert z_{t+h}-\widehat z^{(b)}_{t+h}\rVert_2^2}.
\end{equation}
Positive skill means that DMD improves on that baseline. Report variance-weighted
and equal-PC versions, normalized RMSE, correlation, and the fraction of
non-finite or explosive forecasts.

Required baselines are:

\begin{enumerate}
  \item training mean;
  \item persistence, $\widehat z_{t+h}=z_t$; and
  \item an independent ridge AR model for each PC at the same lag and horizon.
\end{enumerate}

Beating persistence alone is insufficient for the smoothed signal. Full-DMD gain
over diagonal AR is the evidence for cross-component population dynamics. Its
absence does not make stable real relaxation spectra meaningless, but it forbids
a claim of coordinated/rotational modes.

\subsection{Numerical and spectral diagnostics}

For every fit, retain the effective rank, ridge scale, normal-matrix condition
number, residual, spectral radius, non-finite flag, and forecast blow-up flag.
Plot the discrete eigenvalue plane directly. Following Pachitariu et al., report
the number of rotations per tenfold attenuation for sufficiently nonzero complex
eigenvalues, but do not interpret a frequency with fewer than three cycles in the
window or a branch-ambiguous logarithm.

\subsection{Bootstrap and neuron-subset reproducibility}

With the PCA basis fixed, perform 100 moving-block resamples per evaluation
window. Block length is at least 15 native frames (about 2 s) and may increase
with the empirical autocorrelation time. Snapshot pairs are formed only inside
blocks. Group a conjugate pair as one real invariant subspace.

For a development-only reference operator with $q\geq8$, select $k=4$ real
invariant dimensions by mean training modal-coefficient energy, never splitting
a conjugate pair. Match every evaluation and bootstrap fit back to those fixed
reference groups. Compare a bootstrap projector $P^{(b)}$ with its reference
$P$ using
\begin{equation}
  S_{\mathrm{sub}}=
  \frac{1}{k}\tr\!\left(PP^{(b)}\right)
  =\frac{1}{k}\sum_{j=1}^{k}\cos^2\theta_j.
\end{equation}
Also report matched discrete-eigenvalue spread and forecast uncertainty. Repeat
the chosen arm on two deterministic, equal-size neuron subsets to show whether
the conclusion depends on a particular population sample.

Use normalized projector distance
$d_P=\sqrt{\max(0,1-S_{\mathrm{sub}})}$. If fewer than four nontrivial matched
dimensions exist, or if $k=q$, subspace stability is uninformative and the
corresponding gate is not eligible to pass.

To determine whether future tracking has adequate resolution, compute
\begin{equation}
 R_{\mathrm{track}}=
 \frac{\operatorname{median}(\text{within-window bootstrap projector distance})}
      {\operatorname{median}(\text{between-window same-label projector distance})}.
\end{equation}
If estimation noise is not clearly smaller than ordinary between-window change,
a visually smooth Grassmann trajectory would not be trustworthy.

\section{Positive controls and negative controls}

\subsection{Lightweight known-system recovery}

Run 20 deterministic seeds for each of two stable six-dimensional systems: one
with real relaxation modes and one with a known conjugate rotational pair.
Project each latent system into approximately 500 neuronal observables and test
four observation levels, following the causal ordering of the precedent's
observation simulation where applicable:

\begin{enumerate}
  \item continuous Gaussian observations (software upper bound);
  \item rectified latent rates followed by sparse Poisson/event observations,
  matched to the empirical zero fraction;
  \item the same events after a 15-frame Gaussian smoothing approximation; and
  \item a 0.25-s exponential calcium convolution with matched shot noise.
\end{enumerate}

This is intentionally smaller than the end-to-end simulator in the parent plan.
OASIS re-estimation and switching/mode-birth simulations belong to expanded
validation if this first screen passes. The lightweight test must nevertheless
show when sparsity or smoothing changes the recovered eigenspectrum.

\subsection{Empirical negative controls}

For each evaluation window, generate 100 neuronwise circular-shift controls
inside the same continuous, constant-label bout. They preserve each neuron's
marginal distribution and autocorrelation approximately while disrupting
synchronous population relationships. Also simulate independent stationary
event/calcium traces with empirical rates and the same smoothing.

DMD may predict these controls because marginal autocorrelation is genuine. The
null expectation is instead that full-DMD gain over diagonal AR and reproducible
rotational/coupled subspaces disappear. A smoothed trace is not validated merely
because it is easy to forecast.

\section{Predeclared screening gates}

These are feasibility thresholds, not biological hypothesis-test thresholds.
They are frozen before running the pilot and reported even when they fail.

\begin{gatebox}[Gate 1 --- geometry and numerical health]
There must be zero snapshot pairs across an acquisition boundary, raw-label
change, or bootstrap-block join. At least 95\% of evaluation fits must produce
finite eigenvalues and finite forecasts with complete provenance.
\end{gatebox}

\begin{gatebox}[Gate 2 --- known-system recovery]
Across the sparse/calcium-like synthetic observations, at least 80\% of seeds
must correctly distinguish real-relaxation from rotational systems; median
$S_{\mathrm{sub}}$ must be at least 0.75; identifiable frequency/decay error must
be no more than 25\%; and held-out forecast skill must be positive.
\end{gatebox}

\begin{gatebox}[Gate 3 --- empirical predictive validity]
At least one arm must have positive median skill over persistence near 1 s, with
a moving-block-bootstrap lower confidence bound above zero in at least three of
the four recording--state cells. It must not show systematic forecast failure at
about 2 s, and fewer than 5\% of fits may be non-finite or explosive.
\end{gatebox}

\begin{gatebox}[Gate 4 --- empirical reproducibility]
Median $S_{\mathrm{sub}}$ must be at least 0.80 in at least three of the four
recording--state cells, and the qualitative spectral class must agree across the
two deterministic neuron subsets.
\end{gatebox}

\begin{gatebox}[Gate 5 --- tracking resolution]
$R_{\mathrm{track}}<0.5$ in both development recordings. If a 300-frame arm
fails but its prespecified 1,500-frame version passes, the outcome is REVISE the
window rather than PASS the original sliding-window design.
\end{gatebox}

\begin{riskbox}[Coordinated-mode interpretation qualifier]
A claim of multivariate coupling or a reproducible rotational mode additionally
requires full-DMD gain over diagonal AR, or the relevant rotational statistic,
to exceed the 95th percentile of the circular-shift null. Failure of this
qualifier permits only a conditional statement about stable marginal/relaxation
spectra; it does not satisfy a coordinated-mode claim.
\end{riskbox}

\section{Decision rule and the review pause}

\begin{longtable}{@{}P{0.17\textwidth}P{0.48\textwidth}P{0.27\textwidth}@{}}
\toprule
Decision & Evidence pattern & Authorized next step after review\\
\midrule
\endfirsthead
\toprule
Decision & Evidence pattern & Authorized next step after review\\
\midrule
\endhead
PASS & At least one arm passes Gates 1--5. Any coordinated-mode claim also
passes the null qualifier. & Freeze that candidate and propose expanded
validation on all recordings. Do not start tracking automatically.\\
CONDITIONAL & Prediction and stability pass, but diagonal/null evidence is weak,
or only a 1,500-frame/$\Delta F/F$ configuration passes. & Revise signal, window,
or interpretation; limit the endpoint to reproducible decay/eigenspectrum or
whole-subspace summaries.\\
FAIL & Synthetic event recovery fails, no empirical arm beats persistence, or
subspaces are dominated by estimation/null variation. & Stop intensive DMD;
report the identifiability bound and retain static rate/covariance or a
calcium-aware latent-model alternative.\\
\bottomrule
\end{longtable}

The implementation phase ends after generating the artifacts below. The user
reviews the evidence and explicitly decides whether to expand, revise, or stop.

\section{Required implementation artifacts}

All future DMD-specific files remain inside \texttt{DMD/}:

\begin{itemize}
  \item thin entry point: \texttt{DMD/scripts/00\_initial\_dmd\_validation.py};
  \item reusable implementation and tests: \texttt{DMD/src/} and
  \texttt{DMD/tests/};
  \item immutable run output: \texttt{DMD/results/03\_initial\_validation/}; and
  \item result report: \texttt{DMD/documents/04\_dmd\_initial\_validation\_results.pdf}.
\end{itemize}

The result directory must include a machine-readable configuration, software
versions and seeds, selected-window manifest, per-fit metrics, bootstrap metrics,
null distributions, and explicit failure records. No DMD implementation will be
placed in the repository's tutorial \texttt{scripts/}, \texttt{src/funcnet}, or
general \texttt{results/} directories.

Required report figures, in order, are:

\begin{enumerate}
  \item signal distributions, zero fraction, autocorrelation, and representative
  all-neuron/PC traces for each arm;
  \item the Pachitariu-style discrete-eigenvalue plane and rotation statistic;
  \item forecast skill versus physical horizon, arm, lag, and rank, with both
  persistence and diagonal-AR baselines;
  \item representative held-out PC predictions;
  \item bootstrap eigenvalue clouds and invariant-subspace stability;
  \item known-system recovery across observation levels;
  \item observed versus circular-shift/null multivariate evidence; and
  \item a traffic-light table showing every gate, exception, and final decision.
\end{enumerate}

\section{Relationship to the parent analysis plan}

This document is a prospective pilot supplement, not a silent rewrite of
\path{02_dmd_brain_state_analysis_plan.pdf}. The parent document currently
names native \texttt{spike\_smoothed}, unscaled centering, and a one-frame lag as
primary choices. Those choices are now treated as unresolved until this pilot is
reviewed. Raut et al. remain a benchmark for continuous widefield-$\Delta F/F$
DMD, not the precedent for DMD on sparse per-neuron events \citep{raut2025}.

If the pilot passes, its locked signal, scaling, lag, rank, estimator, and window
will be written as a versioned amendment to the parent plan before state labels
are used for biological comparison. If the pilot fails, the negative result is
the scientifically correct stopping point.

\clearpage
\sloppy
\bibliographystyle{unsrtnat}
\bibliography{dmd_tracking_references}

\end{document}
